\section{Related Work and Context}

\subsection{Cloud-Native Orchestration for the Edge}
Kubernetes-native edge orchestration approaches commonly extend cluster semantics through custom resources, controllers, and the operator pattern~\cite{kubernetes_operators}. In most cases, however, the managed endpoints are Linux-capable edge nodes or gateway-class systems able to host containers, agents, or complete Kubernetes distributions. This assumption simplifies orchestration because the control plane can rely on relatively rich host capabilities and standard networking behavior. RETROSPECT targets a more constrained setting in which endpoint devices are embedded-class targets and therefore cannot be treated as ordinary worker nodes. The platform preserves the declarative orchestration mindset of Kubernetes while explicitly decoupling it from the assumption that every managed endpoint can run cloud-native software stacks natively.

\subsection{Embedded RTOS and WASM Runtime Integration}
Zephyr provides a production-grade RTOS with a wide hardware abstraction layer, networking support, and security features suitable for connected embedded systems~\cite{zephyr}. WAMR complements this landscape by offering a lightweight runtime capable of executing WebAssembly modules in constrained environments~\cite{wamr}. Prior work and engineering demonstrations have established the feasibility of running WASM beyond the browser and even beyond traditional servers~\cite{haas_wasm}. RETROSPECT builds on this enabling substrate but shifts the emphasis from local runtime feasibility to platform-level orchestration. In other words, the question is not simply whether a microcontroller can execute a WASM artifact, but whether a cloud control plane can securely assign, observe, and supervise that execution through a coherent operational model.

\subsection{Secure Communication and Binary Encoding}
TLS 1.3 provides the baseline for secure transport with modern cryptographic negotiation and strong protection against passive and active attacks on the communication channel~\cite{rfc8446}. CBOR provides a compact binary encoding well suited to constrained networks and low-overhead message exchange~\cite{rfc8949}. Together, these technologies are a sensible foundation for southbound communication in embedded orchestration scenarios. Nevertheless, transport security and compact serialization do not automatically define resource ownership, identity binding, or state propagation semantics. RETROSPECT therefore contributes not a new secure transport, but an integrated use of TLS and CBOR inside a gateway-mediated life-cycle workflow aligned with Kubernetes-style status management.

\subsection{Emulation-Driven Validation and Reproducibility}
Renode has emerged as a valuable framework for embedded development and test automation because it enables hardware-faithful emulation and repeatable multi-node scenarios~\cite{renode}. In experimental systems research, such tooling is particularly important when physical hardware inventories are limited, when repeated regression testing is required, or when peer reviewers need a realistic description of the validation path. RETROSPECT leverages emulation not as a substitute for all hardware experiments, but as an enabler of reproducible control-plane and protocol validation. This is aligned with the paper's scope: the emphasis is on the correctness and operational coherence of the orchestration stack rather than on cycle-level microarchitectural benchmarking.

\subsection{Positioning of the Present Contribution}
Relative to the state of practice, RETROSPECT occupies a distinct systems integration space. It does not merely present a lightweight runtime, a secure device protocol, or a Kubernetes extension in isolation. Instead, it documents how a cloud control plane, a fog-layer gateway, and constrained edge endpoints can be combined into an operationally meaningful whole. This positioning is important because the most difficult engineering problems arise at the boundaries between these layers: who owns reported state, how identity is verified, how desired and observed state are synchronized, and how deployment experiments can be reproduced. These concerns define the paper's contribution more accurately than any individual implementation detail.
